\section{Modélisation}\label{sec:methodes}

\subsection{Protocole d'évaluation}\label{sec:protocole}

Tous les modèles suivent le même protocole, pour être comparables.

\subsubsection{Variables à prédire}\label{sec:qualite}

\paragraph{Pourquoi ces grandeurs.} La base ne dit pas si une soudure est « bonne ».
Nous mesurons sa qualité par les propriétés que vérifient les normes de soudage : la
\emph{résistance} et la \emph{ductilité} (essai de traction), et la \emph{ténacité},
résistance aux chocs (essai Charpy). L'absence de défauts ou la tenue en fatigue ne sont
pas dans la base.

\paragraph{Ce que les données en disent.}
\begin{itemize}
  \item Limite d'élasticité et résistance à la traction sont presque redondantes
    (corrélation 0,93 sur 714 lignes).
  \item La résistance s'oppose à la ductilité (corrélations de $-0{,}57$ à $-0{,}78$) :
    une bonne soudure est un compromis.
  \item Les lignes Charpy sont de deux types : 405 énergies mesurées, 474 fixées à
    \SI{100}{\joule} ou \SI{28}{\joule} (section~\ref{sec:exploration}).
\end{itemize}

\begin{table}[!htbp]
  \centering\small
  \begin{tabular}{@{}l l c r r@{}}
    \toprule
    \textbf{Qualité} & \textbf{Grandeur} & \textbf{Unité} & \textbf{Lignes} & \textbf{Articles} \\
    \midrule
    Résistance & \textbf{limite d'élasticité} & \si{\mega\pascal} & 780 & 42 \\
               & résistance à la traction     & \si{\mega\pascal} & 738 & 41 \\
    Ductilité  & \textbf{allongement}         & \%                & 700 & 36 \\
               & striction                    & \%                & 705 & 34 \\
    Ténacité   & \textbf{énergie Charpy}      & \si{\joule}       & 879 & 26 \\
    \bottomrule
  \end{tabular}
  \caption{Grandeurs mesurant la qualité. En gras, les trois cibles principales.}
  \label{tab:cibles}
\end{table}

\paragraph{Stratégie.}
\begin{enumerate}
  \item \textbf{Prédire chaque propriété} : une cible principale par qualité (en gras,
    tableau~\ref{tab:cibles}), chacune avec son modèle de régression, entraîné sur les
    lignes où elle est mesurée. Pour le Charpy, la température d'essai est une entrée ;
    c'est une condition de mesure, pas un réglage : les recommandations seront données
    pour une température d'essai fixée.
  \item \textbf{Juger la qualité} : la soudure est bonne si chaque valeur prédite dépasse
    son seuil.
\end{enumerate}
Un modèle direct « bonne / mauvaise » serait moins adapté : les seuils dépendent de
l'usage, un oui/non perd de l'information (avec un seuil de \SI{47}{\joule},
\SI{46}{\joule} et \SI{10}{\joule} seraient tous deux « mauvais »), et résistance et
ténacité sont rarement mesurées sur la même ligne.

Ce choix exclut les modèles de classification, pas les familles de modèles (arbres,
forêts, plus proches voisins existent en régression). Son coût : la régression ne se
concentre pas sur les erreurs proches du seuil. Une classification est donc testée en
complément sur le Charpy, seule propriété avec un seuil commun à tous les aciers
(\SI{27}{\joule} ou \SI{47}{\joule}) ; le seuil de résistance dépend de la classe d'acier
(section~\ref{sec:modeles}).

\outrouver{\code{TARGETS} et \code{get_xy} dans \code{preprocessing/ml_dataset.py} ;
notebook, étape 6 et « Analyse descriptive › Liens entre les grandeurs à prédire ».}

\afaire{\item Valider en équipe les trois cibles principales.
  \item Trouver les seuils de « bonne soudure » et leurs sources (normes ISO 2560,
    ISO 15614-1, AWS A5.1 à vérifier ; seuils Charpy usuels de \SI{27}{\joule} ou
    \SI{47}{\joule} à une température donnée).}

\subsubsection{Validation croisée}

Un modèle est jugé sur des données qu'il n'a pas vues : la \emph{validation croisée}
découpe les données en 5 parts, entraîne sur 4, teste sur la cinquième, cinq fois, puis
moyenne les scores. Les lignes d'un même cordon ou d'un même article étant très proches,
un découpage au hasard serait trop optimiste. Les lignes sont donc regroupées
\textbf{par article} (\code{GroupKFold}) : le score mesure la capacité à prédire les
soudures d'une étude inconnue.

\outrouver{fonction \code{weld_source} dans \code{preprocessing/ml_dataset.py} ;
vérification de l'absence de fuite : notebook, étape 6.}

\afaire{\item Validation croisée imbriquée : un niveau extérieur (par article) pour mesurer
    la performance, un niveau intérieur (par article) pour régler les hyperparamètres.
  \item Ne pas utiliser l'erreur OOB des forêts : elle ignore les articles et serait trop
    optimiste.}

\subsubsection{Mesures d'erreur}

\afaire{\item Justifier les mesures : RMSE et MAE, dans l'unité de la cible
    (\si{\mega\pascal}, \%, \si{\joule}), et $R^2$ ; moyenne $\pm$ écart-type sur les plis.
  \item Choix final par la règle de l'écart-type (cours 2) : parmi les modèles presque
    aussi bons que le meilleur, garder le plus simple.
  \item Coder le protocole commun (modèle + cible $\rightarrow$ scores), à faire en
    premier --- \textbf{Salomé (proposé)}.}

\subsection{Entraînement et évaluation des modèles}\label{sec:modeles}

Chaque famille de modèles est appliquée aux trois cibles principales.

\subsubsection{Modèles de référence et modèles linéaires}

\afaire{\item Référence : prédire la moyenne (tout modèle doit faire mieux).
  \item Régression linéaire, sélection de variables pas à pas, Ridge, Lasso ; lecture des
    coefficients --- \textbf{Salomé (proposé)}.}

\subsubsection{Réduction de dimension}

\afaire{\item Régression sur les composantes de l'ACP et PLS ; nombre de composantes réglé
    par validation croisée.
  \item Charpy : modéliser les deux types de lignes ensemble ou séparément ? Tester les
    deux --- \textbf{?}}

\subsubsection{Arbres et méthodes d'ensemble}

\afaire{\item Arbre CART, bagging, forêt aléatoire, extra-trees ; hyperparamètres
    (profondeur, taille minimale des feuilles, nombre de variables tirées).
  \item Importance des variables par permutation ; attention aux variables corrélées
    (Cr, Mo, traitement), dont l'importance est surestimée --- \textbf{?}
  \item Classification complémentaire « Charpy $\geq$ \SI{47}{\joule} à la température
    d'essai ? » : comparer un arbre ou une forêt de classification au même modèle en
    régression suivi du seuil, sur les mêmes plis.
    \begin{itemize}
      \item Seulement les 405 énergies mesurées (13 articles) : sur les lignes fixées,
        l'étiquette est imposée par la convention (\SI{100}{\joule} toujours « bonne »,
        \SI{28}{\joule} toujours « mauvaise »).
      \item Classes déséquilibrées (69\,\% de « bonnes » à \SI{47}{\joule}, 84\,\% à
        \SI{27}{\joule}) : précision équilibrée, F1, matrice de confusion.
      \item Piste facultative : classes de résistance de la norme des électrodes
        (ISO 2560, seuils à vérifier) pour la traction.
    \end{itemize}}

\subsubsection{Apprentissage semi-supervisé}

\afaire{\item Bibliographie et au moins une méthode (auto-apprentissage, co-régression
    COREG) exploitant les lignes où la cible n'est pas mesurée ; justifier son
    intérêt --- \textbf{?}}

\subsection{Conclusion de la modélisation}\label{sec:resultats}

\afaire{\item Tableau comparatif des modèles par cible (moyenne $\pm$ écart-type sur les plis).
  \item Figures prédit / réel.
  \item Poids des indicateurs « valeur manquante » (section~\ref{sec:prep-manquantes}) :
    comparer avec et sans.
  \item Modèle retenu pour chaque cible et limites (peu d'articles, extrapolation).
  \item Chacun présente sa famille de modèles aux autres avant le partiel.}

% Gabarit de tableau à remplir
% \begin{table}[!htbp]
%   \centering\small
%   \begin{tabular}{@{}lccc@{}}
%     \toprule
%     \textbf{Modèle} & \textbf{RMSE (MPa)} & \textbf{MAE (MPa)} & $\boldsymbol{R^2}$ \\
%     \midrule
%     Ridge  & & & \\
%     \bottomrule
%   \end{tabular}
%   \caption{Limite d'élasticité, validation croisée groupée par article.}
% \end{table}
